% =====================================================================
%  CHAPTER 5: 儿童少年身体发育
% =====================================================================
\chapter{儿童少年身体发育}
\setcounter{chapter}{5}
\chapterintro{
掌握儿童体格发育的阶段性变化（四个阶段）；掌握第一和第二生长突增期的特点；掌握体成分发育变化；掌握体能（生理功能和运动能力）的发育规律；熟悉体格发育的性别差异；了解身体比例和体型的发育变化。
}

% ----- 5.1 体格发育 -----
\section{体格发育}

\subsection{体格发育的概念}

\begin{defbox}
\textbf{体格发育（Physical Growth）}：指身体形态的不断增大（身高、体重、器官（除淋巴组织外）大小），是一个形态不断增加和重塑的过程，贯穿婴幼儿期、儿童期和青春期。体格发育水平反映生长发育状况，为评价营养、运动、健康水平以及预防干预的可及性提供重要依据。

\textbf{体格发育特点：}形态性、动态性、阶段性、连续性。
\end{defbox}

\subsection{体格发育的四个阶段}

\begin{keybox}
\textbf{基于身高发育速度曲线的四个阶段（Tanner JM, 1966）：}

\begin{enumerate}[leftmargin=*]
    \item \imp{第一阶段：第一生长突增期（First Stage of Growth Spurt）}\\
    从孕期开始至出生后1--2岁末。
    \begin{itemize}
        \item \textbf{身高}：孕期4--6个月期间增长量约占胎儿期总增长量的50\%（约25cm），出生后第1年增长约25cm
        \item \textbf{体重}：孕期7--9个月期间增长量约占胎儿期总增长量的70\%
    \end{itemize}

    \item \imp{第二阶段：相对稳定期}\\
    2岁至青春期发育前。
    \begin{itemize}
        \item 身高：约5--7cm/年
        \item 体重：约2--3kg/年
    \end{itemize}

    \item \imp{第三阶段：第二生长突增期（Adolescent Growth Spurt）}\\
    青春期阶段。
    \begin{itemize}
        \item 身高由每年增长5--10cm开始，逐渐达突增高峰（PHV），一般为10--14cm/年，男生＞女生
        \item 体重每年增加4--7kg，突增高锋时一年可增8--10kg
    \end{itemize}

    \item \imp{第四阶段：生长停止期}\\
    青春期后期开始，身高基本停止增长，体重也一般停止增长或仅少量增加。
\end{enumerate}
\end{keybox}

\subsection{身体比例的变化}

\begin{defbox}
\textbf{反映身体充实度的指标（横向指标）：}

\begin{enumerate}[leftmargin=*]
    \item \imp{身高胸围指数}（胸围/身高）和\imp{身高体重指数}（体重/身高）：\\
    只能衡量身体局部的充实度，上部量和下部量的实际值重复性差，身高代替各部分发育的差异，在青春期表现更明显

    \item \imp{BMI（身体质量指数）}：\\
    能较准确反映不同群体、个体的胖瘦变化。青春期发育使体重快速增加，BMI高百分位数的变化趋势是制定肥胖防控策略和措施的重要依据
\end{enumerate}

\textbf{反映身体比例的指标（纵向指标）：}

\begin{itemize}[leftmargin=*]
    \item \imp{身高坐高指数} = 坐高/身高
    \item \imp{下肢长指数I} = （身高 - 坐高）/身高 × 100\%
    \item \imp{下肢长指数II} = （身高 - 坐高）/坐高 × 100\%
\end{itemize}
\end{defbox}

\subsection{体型的发育变化}

\begin{defbox}
\textbf{体型（Somatotype）}：是对人体形态的总体描述和评分，包括人体整体和各部位大小、身体形态和功能的综合。由遗传因素决定，也受机体对环境适应和与环境因素的综合影响。

\textbf{Sheldon体型分类法（1940年代）：}
\begin{enumerate}[leftmargin=*]
    \item \imp{内胚层型（Endomorphy）}：身体圆润，消化器官发达
    \item \imp{中胚层型（Mesomorphy）}：身体健壮，骨骼和肌肉发达
    \item \imp{外胚层型（Ectomorphy）}：身体瘦长，神经系统和感觉器官发达
\end{enumerate}

目前更多应用\textbf{Heath-Carter体型分类法}，主要用于运动员体型研究。
\end{defbox}

% ----- 5.2 体能发育 -----
\section{体能发育}

\subsection{体能的概念与分类}

\begin{defbox}
\textbf{体能（Physical Fitness）}：又称体适能，指个体具备胜任日常生活和学习而不感到疲劳，同时有余力享受休闲活动、抵抗应激和突发紧急状况的身体能力。1985年Caspersen CJ等将体能分为\imp{健康相关体能}和\imp{运动技能体能}两类。

\textbf{健康相关体能}：与健康水平密切相关，包括心肺耐力、肌肉力量与耐力、柔韧性、身体成分等。

\textbf{运动技能体能}：与运动表现相关，包括灵敏性、平衡、协调、速度、爆发力、反应时等。
\end{defbox}

\subsection{生理功能发育}

\begin{keybox}
\textbf{（一）心血管功能发育}

\begin{itemize}[leftmargin=*]
    \item \imp{心率}：随年龄增长逐渐减慢，新生儿约120--140次/分，学龄儿童约80--90次/分，成人约60--100次/分
    \item \imp{血压}：随年龄增长逐渐升高。儿童期收缩压约90--110mmHg，舒张压约60--70mmHg
    \item \imp{每搏输出量和心输出量}：随年龄增长而增加
\end{itemize}

\textbf{（二）呼吸功能发育}

\begin{itemize}[leftmargin=*]
    \item \imp{呼吸频率}：新生儿约40--44次/分，1岁约30次/分，学龄儿童约20--22次/分，成人约12--16次/分
    \item \imp{肺活量}：随年龄增长显著增加，男生一般＞女生（青春期后差异更明显）
    \item \imp{最大摄氧量（VO$_2$max）}：反映全身有氧耐力水平，青春前期男女差别不大，青春期后男生明显高于女生
\end{itemize}

\textbf{（三）综合功能发育指标}

\begin{itemize}[leftmargin=*]
    \item \imp{每千克体重肺活量指数}：肺活量（ml）/ 体重（kg）
    \item \imp{Blanche心功能指数（BI）}
\end{itemize}
\end{keybox}

\subsection{运动能力发育}

\begin{defbox}
\textbf{运动能力的主要方面：}

\begin{enumerate}[leftmargin=*]
    \item \imp{力量（Strength）}：如俯卧撑、引体向上、握力、立定跳远、仰卧起坐、掷铅球等
    \item \imp{耐力（Endurance）}：20m折返跑、50m×8往返跑（小学）、1000m跑（中学男生）、800m跑（中学女生）、台阶运动试验
    \item \imp{速度（Speed）}：短跑、游泳、滑雪等
    \item \imp{灵敏性（Agility）}：10m×4往返跑、跳绳、乒乓球、平衡木等
    \item \imp{柔韧性（Flexibility）}：坐位体前屈、立位体前屈
    \item \imp{平衡（Balance）}
    \item \imp{协调性（Coordination）}
    \item \imp{反应时（Reaction Time）}
\end{enumerate}

\textbf{运动能力发育的一般规律：}
\begin{itemize}[leftmargin=*]
    \item 男生运动能力普遍优于女生（青春期后差异更明显）
    \item 各种运动能力有各自的敏感发展期
    \item 青春期前后是运动能力快速发展时期
\end{itemize}
\end{defbox}

% ----- 5.3 体成分发育 -----
\section{体成分发育}

\subsection{体成分的概念}

\begin{defbox}
\textbf{体成分（Body Composition）}：指人体总重量中脂肪组织和非脂肪组织（去脂体重，包括骨骼、肌肉、水分等）所占的百分比。体成分反映个体的营养状况和身体发育水平。

\textbf{常用测定方法：}
\begin{enumerate}[leftmargin=*]
    \item \imp{皮褶厚度（Skinfold Thickness）}：测定肱三头肌、肩胛下角等部位，间接估计体脂率
    \item \imp{生物电阻抗法（Bioelectrical Impedance Analysis, BIA）}：通过微弱电流通过身体的阻抗估计体成分
    \item \imp{双能X射线吸收法（DEXA）}：金标准，精确测定脂肪、肌肉和骨密度
\end{enumerate}
\end{defbox}

\subsection{体成分的年龄变化}

\begin{keybox}
\textbf{不同发育阶段体成分的变化特点：}

\begin{enumerate}[leftmargin=*]
    \item \imp{婴幼儿期}：体脂率较高（"婴儿肥"），随后逐渐下降
    \item \imp{学龄前期至学龄期}：体脂率进一步下降（约5--7岁为"体脂反弹期"），瘦体重比例增加
    \item \imp{青春期}：性别差异显著
    \begin{itemize}
        \item \textbf{男性}：体脂率下降或保持较低水平，肌肉和骨骼量显著增加（雄激素作用）
        \item \textbf{女性}：体脂率升高（雌激素促进脂肪沉积），以臀部和大腿为主，为生育做准备
    \end{itemize}
    \item \imp{体脂反弹（Adiposity Rebound）}：儿童期体脂率由降转升的转折点，约在5--7岁。\keyword{体脂反弹越早，儿童期肥胖的风险越高}
\end{enumerate}
\end{keybox}

\subsection{体成分的性别差异}

\begin{table}[H]
\centering
\caption{体成分的性别差异（青春期后）}
\begin{tabularx}{\textwidth}{l>{\raggedright\arraybackslash}p{3.5cm}>{\raggedright\arraybackslash}X}
\hline
\textbf{指标} & \textbf{男性} & \textbf{女性} \\
\hline
体脂率 & 较低（约15--18\%）& 较高（约22--28\%）\\
瘦体重（肌肉量）& 较高 & 较低 \\
骨密度 & 较高 & 较低（绝经后加速流失）\\
体脂分布 & 腹部（中心性/苹果型）& 臀部/大腿（周围性/梨型）\\
\hline
\end{tabularx}
\end{table}

% ----- 5.4 性别发育 -----
\section{性别发育}

\begin{defbox}
\textbf{性别发育（Sex Development）}：包括性别的生物学基础（染色体性别、性腺性别、表型性别）和性别认同的心理学方面。

\begin{enumerate}[leftmargin=*]
    \item \imp{染色体性别}：XX（女性）或XY（男性），在受精时决定
    \item \imp{性腺性别}：卵巢或睾丸的分化（受SRY基因调控）
    \item \imp{表型性别}：内外生殖器和第二性征的发育
    \item \imp{性别认同}：个体对自己性别的心理认知
\end{enumerate}

男女在体格、体成分、体能发育等方面存在显著的生物学差异，这些差异在青春期后愈发明显。
\end{defbox}

% ----- 复习题 -----
\section{复习题}

\begin{enumerate}[leftmargin=*, itemsep=8pt]
    \item \textbf{【名词解释】}体格发育；体能；体成分；体脂反弹
    \item \textbf{【名词解释】}BMI；皮褶厚度；最大摄氧量（VO$_2$max）
    \item \textbf{【名词解释】}Sheldon体型分类；身高坐高指数
    \item \textbf{【简答题】}简述儿童体格发育的四个阶段及各自特点。
    \item \textbf{【简答题】}简述儿童体能发育的主要方面及评估指标。
    \item \textbf{【简答题】}试述儿童体成分发育的年龄变化及性别差异。
    \item \textbf{【简答题】}简述体脂反弹的概念及其公共卫生学意义。
\end{enumerate}

\subsection*{参考答案要点}

\begin{enumerate}[leftmargin=*, itemsep=5pt]
    \item \textbf{体格发育}：身体形态不断增大和重塑的过程。\textbf{体能}：个体胜任日常活动而不疲劳且有余力的身体能力，分为健康相关和运动技能两类。\textbf{体成分}：脂肪组织与去脂体重在总体重中的构成比例。\textbf{体脂反弹}：儿童期体脂率由降转升的转折点，约5--7岁。

    \item \textbf{BMI}：体重(kg)/身高$^2$(m$^2$)，反映体重是否在健康范围。\textbf{皮褶厚度}：用皮褶计测量特定部位皮下脂肪厚度，间接估计体脂率。\textbf{VO$_2$max}：最大摄氧量，反映全身有氧耐力。

    \item \textbf{Sheldon分类}：将体型分为内胚层型（圆胖）、中胚层型（健壮）、外胚层型（瘦长）。\textbf{身高坐高指数}：坐高÷身高，反映身体上下部比例。

    \item 四个阶段：(1)第一生长突增期（孕期至1--2岁末，身高增约25cm/年）；(2)相对稳定期（2岁至青春期前，5--7cm/年）；(3)第二生长突增期（青春期，PHV达10--14cm/年，男＞女）；(4)生长停止期（青春期后期）。

    \item 生理功能：心血管（心率、血压）、呼吸（肺活量、VO$_2$max）；运动能力：力量、耐力、速度、灵敏性、柔韧性、平衡、协调、反应时。

    \item 婴幼儿期体脂率高→学龄前期下降→体脂反弹（5--7岁）→青春期性别差异显著（男肌肉增多体脂下降，女体脂上升）。体脂反弹越早，肥胖风险越高。

    \item 体脂反弹约5--7岁，是儿童期体脂率由降转升的转折点。越早发生，儿童期肥胖风险越高。公共卫生意义：该时期是预防肥胖的关键窗口期，应加强膳食和运动干预。
\end{enumerate}

\newpage
